% !TeX root = ../tesis.tex

Con la maquinaria de \Mithril ya descrita, podemos volver a la implementación del \zkBridge y leerla con más precisión. El \StakeDistributionUTxO en \ChainB es la persistencia \Onchain de una historia local de certificados de \Mithril. Su fase de \textit{bootstrap} ancla el sistema en un \GenesisCertificate; sus actualizaciones recurrentes consumen un certificado padre $\Certificate_{p,n}$ o $\FirstCertificate_n$ y publican uno nuevo, preservando el mismo NFT y, por lo tanto, el mismo punto de verdad para la capa de actualización del \zkBridge. El \StakeDistributionUTxO mantiene disponible una versión local de $\CertificateChain$ y del par $(\SnarkAuthReg, \SnarkAVK)$ que el \zkBridge acepta como vigente. Así, la autenticidad del estado depende de la verificabilidad de la \CertificateChain de \Mithril, y no de un operador particular.

El flujo de \MintTx también se aclara al mirar \Mithril desde esta perspectiva. No hace falta validar el consenso de \Cardano. Necesitamos validar un \TransactionSnapshotCertificate, un argumento de que ese certificado es válido respecto del estado de \StakeDistribution disponible como \textit{reference input}, y un argumento de pertenencia de la \LockTx dentro del \TransactionSnapshot certificado.

Esto también explica la decisión, ya anticipada en la primera etapa de implementación, de no rehashear certificados completos \Onchain. Una vez que se entiende que un certificado ya viene ligado a un \ProtocolMsg, a una \SnarkAVK, a una \CertificateChain y a un tipo de artefacto, recomputar su hash entero dentro del validador deja de ser la única manera de obtener \textit{binding}. La implementación final es más económica.

En el caso de nuestro \ZkBridgeOperator, esta interfaz aparece concretamente en la separación entre artefactos persistentes y evidencia efímera. Los certificados sincronizados y el estado de \StakeDistribution viven en el espacio de estado del protocolo y en el almacenamiento del operador. Los argumentos de pertenencia de transacciones, los \textit{witnesses} normalizados y los datos útiles para \Groth viven sólo durante la operación que los necesita.

Una vez entendido qué certifica \Mithril, cómo lo encadena y cómo lo consume el operador del \zkBridge, ya podemos analizar qué parte de esa información debe formar parte de $\PublicInputs$/$\PrivateInputs$ respectivamente, cuáles chequeos deben hacerse \Onchain y cuáles \Offchain, y cómo impacta todo esto en costos reales de memoria y CPU. El capítulo siguiente retoma esta base.
